% HD-R09. GPT-6 Astra Pro /root, inherited effort, 21 November 2026.
\subsection{HD-R09: line-power sums cannot supply a weakly semiregular imaginary secant}
\label{sec:line-power-secant-obstruction}

Let $X$ be a complex abelian variety of even dimension $m=2n\ge4$.
Let $H,B\in H^{1,1}(X,\mathbf R)$, with $H$ nondegenerate.
Consider actual finite direct sums
\[
 F=\bigoplus_{j=1}^N L_j[a_j],\quad
 L_j\in\operatorname{Pic}(X),\quad a_j\in\mathbf Z,
 \qquad c_1(L_j)=B+t_jH,\quad t_j\in\mathbf R.              \tag{R09.1}
\]
Every such object with whole character
\[
       \operatorname{ch}(F)=e^B
              \bigl(\alpha e^{sH}+\overline\alpha e^{-sH}\bigr),
       \qquad s=i\sqrt d,\quad d>0,\quad\alpha\ne0,          \tag{R09.2}
\]
fails weak semiregularity for every actual coherent action of a
faithful finite point subgroup $G\subset X\times\widehat X$.
This includes arbitrary mixing linearizations, repeated line bundles,
flat twists and integral shifts. The proof gives an invariant
degree-two trace kernel of dimension at least $\binom m2$.
No degree-$(m-2)$ trace injectivity is used. In particular the
index-three source remaining in HD-R08 cannot be supplied in
dimension six by the actual objects(R09.1).

\paragraph{1. Actual invariant orbit summands.}
We use the actual idempotent construction of HD-N13, and record
its maps here because the finite action need not respect the
chosen direct-sum coordinates. Each $L_j[a_j]$ has endomorphism
algebra $\mathbf C$. Krull--Schmidt decomposition groups the
repeated isomorphism types into matrix blocks in
$E/J$, where $E=\operatorname{End}(F)$ and $J$ is its
nilpotent Jacobson radical. Translation and tensoring by a flat
line bundle preserve the shift and the first Chern class.
Thus the action permutes these types, and every orbit has
one shift and one value $t_\nu$ in(R09.1).

For each orbit, its central projector in $E/J$ is invariant.
Averaging over $G$ gives the surjection
$E^G\to(E/J)^G$. Idempotents lift through the nilpotent
kernel $J\cap E^G$ by the explicit Newton step
\[
                    e\longmapsto e-(2e-1)^{-1}(e^2-e).
                                                               \tag{R09.3}
\]
Indeed $(2e-1)^2=1+4(e^2-e)$ is invertible, and the new
error is the square of the old error times a commuting unit.
Nilpotence makes the iteration terminate. Successive
complementary corners give mutually orthogonal invariant
projectors summing to one. Their actual splittings in the
finite equivariant category give
\[
                 F\simeq\bigoplus_{\nu=1}^kF_\nu,
           \qquad\operatorname{ch}(F_\nu)=w_\nu e^{B+t_\nu H},
           \quad w_\nu\in\mathbf Z\setminus\{0\}.          \tag{R09.4}
\]
These have the original orbit-sum underlying objects. Explicitly,
idempotents $e,f$ with the same reduction modulo $J$ are
intertwined by the unit $u=fe+(1-f)(1-e)$, which satisfies
$ue=fu$ and $u\equiv1\pmod J$. The common shift in each
orbit makes its signed multiplicity $w_\nu$ nonzero.

The invariant scalar maps
\[
 H^2(X,\mathcal O_X)\longrightarrow
            \operatorname{Ext}^2(F_\nu,F_\nu)^G,
              \qquad\beta\longmapsto\beta\operatorname{id}_{F_\nu}
                                                               \tag{R09.5}
\]
are injective: their degree-zero Chern trace component is
$w_\nu\beta$. They are invariant because translations and
flat tensoring act identically on $H^*(\mathcal O_X)$.
The independent diagonal summands in(R09.4) therefore give
\[
                  \dim\operatorname{Ext}^2(F,F)^G
                                      \ge k\binom m2.        \tag{R09.6}
\]

\paragraph{2. The exact number of distinct real exponent nodes.}
Group equal $t_\nu$ in(R09.4), and write $c_t$ for the sum
of their nonzero integer weights. Let $t_1,\ldots,t_r$ be
the distinct nodes for which $c_t\ne0$. Nodes whose total
weight cancels can be omitted from this character calculation;
their actual summands remain in(R09.6), so $k\ge r$.
Since $H$ is nondegenerate, each $H^j$ for $0\le j\le m$
is nonzero. Comparing every homogeneous component in(R09.2)
after multiplication by $e^{-B}$ gives the exact moments
\[
       M_j=\sum_{t}c_t t^j
                =\alpha s^j+\overline\alpha(-s)^j,
                               \quad0\le j\le m.             \tag{R09.7}
\]
At least one of $M_0,M_1$ is nonzero: their simultaneous
vanishing gives $\alpha=0$ by adding and subtracting the two
equations, since $s\ne0$. Thus the node list is nonempty.
The identity $s^2=-d$ gives
\[
              \sum_t c_t(t^2+d)t^j=0,
                              \quad0\le j\le m-2.            \tag{R09.8}
\]
For $r\le m-1$, take the first $r$ equations. Their coefficient
matrix is the Vandermonde matrix $(t_\ell^j)$, with determinant
$\prod_{a<b}(t_b-t_a)\ne0$. Its exact inverse therefore
forces every $c_t(t^2+d)$ to vanish. Each $t$ is real and
$d>0$, so $t^2+d>0$, contradicting the definition of the
nonzero coefficients. Consequently
\[
                               k\ge r\ge m.                 \tag{R09.9}
\]
This argument permits either sign for each integer weight and
does not assume that the virtual character was given by an
effective sum.

\paragraph{3. The actual higher evaluation factors through a smaller space.}
Put $q=m-2$, which is even, and let
$U=H^{0,1}(X)\oplus H^{1,0}(X)^*=HH^1(X)$.
For an actual line bundle with $c_1=B+tH$, the degree-one
characteristic action is the literal map
\[
          \pi_t(\eta,v)=\eta+\iota_v B+t\iota_vH
                                      \quad\text{in }H^{0,1}(X).
                                                               \tag{R09.10}
\]
It follows from its Atiyah class $c_1(L)$; the multiplicative
Hochschild action on the abelian variety gives degree-$q$
evaluation $\bigwedge^q\pi_t$. Its actual target is
$\operatorname{Ext}^q(L,L)=H^q(\mathcal O_X)$.
The exterior Hochschild algebra and its characteristic map
are justified by the group-difference Koszul maps and the
Buchweitz--Flenner locators checked in HD-N19.

The expression $\bigwedge^q\pi_t$ is a polynomial in the
real number $t$ of degree at most $q$. Write its actual
coefficient maps as
\[
   \bigwedge^q\pi_t=\sum_{j=0}^q t^jT_j,
       \qquad T_j:\bigwedge^qU\longrightarrow H^q(\mathcal O_X).
                                                               \tag{R09.11}
\]
The direct-sum evaluation is diagonal on the original line
summands by naturality of the characteristic transformation.
Repeated line bundles receive the identity matrix on their
multiplicity space, and flat twists leave(R09.10) unchanged.
The shift sign in a degree-$q$ characteristic action is
$(-1)^{a_jq}=1$, since $q$ is even. Equivalently the even
action commutes with the canonical shift identifications.
Conjugating by the actual isomorphism to(R09.4) does not
change the dimension of its image.

Thus its image factors through the specified finite-dimensional
space and evaluation-at-node map
\[
 \bigwedge^qU\xrightarrow{(T_0,\ldots,T_q)}
     H^q(\mathcal O_X)^{q+1}
       \xrightarrow{(z_j)\mapsto(\sum t_\nu^jz_j)_\nu}
               \operatorname{Ext}^q(F,F).                  \tag{R09.12}
\]
In the last arrow, each coordinate multiplies the actual
identity of its orbit summand. Formula(R09.12) therefore
proves the uniform bound
\[
       \dim\operatorname{im}\operatorname{ev}_F^{m-2}
                    \le(m-1)\binom m{m-2}
                     =(m-1)\binom m2.                       \tag{R09.13}
\]
It is an image bound on the entire Hochschild evaluation;
no generation assertion about all self-Ext classes is involved.

\paragraph{4. Invariant Serre adjunction gives the trace kernel.}
For a nonidentity point autoequivalence of $X$, its twisted
Hochschild sector vanishes. A nonzero translation has its
graph disjoint from the diagonal. For zero translation and
a nontrivial flat line bundle $P$, all cohomology of $P$
vanishes, giving the same conclusion. The actual diagonal
kernel and forgetful maps of HD-N12 therefore identify weak
$G$-semiregularity with injectivity of the ambient degree-two
trace on $\operatorname{Ext}^2(F,F)^G$, in this dimension
as well. Translation and flat tensoring preserve the invariant
top holomorphic volume and its Serre pairing. The dual of
that invariant trace is consequently the degree-$(m-2)$
evaluation map into $\operatorname{Ext}^{m-2}(F,F)^G$.
These are the same volume-form and averaging morphisms of
HD-N12 and HD-N19; their degrees here are explicitly $2$
and $m-2$.

The rank of a finite-dimensional map equals the rank of its
dual. Equations (R09.6), (R09.9) and (R09.13), together with
invariant Serre duality, therefore give
\[
 \begin{split}
 \dim\ker\bigl(\sigma_F^2|_{\operatorname{Ext}^2(F,F)^G}\bigr)
 &\ge k\binom m2-(m-1)\binom m2\\
 &\ge\binom m2>0.
 \end{split}                                                 \tag{R09.14}
\]
This proves the asserted obstruction and quantitative kernel
for every coherent action in the stated scope.

\paragraph{Scope in the construction problem.}
For $m=6$ the proven kernel lower bound is15; for $m=8$ it
is28. The argument requires an actual direct sum of line
shifts whose first Chern classes lie on the displayed real
affine line. It does not replace a complex with nonzero
differential by its virtual line-bundle expansion. The exact
maps(R09.3),(R09.10)--(R09.12) retain that object-level
requirement. HD-N13 treats a different fourfold scope with
arbitrary divisor classes; the present proof closes the
single-class line-power route in every even dimension,
including the indefinite sources left by HD-R08. In the
project frame, these $K$-theory character expansions do
not give the required deformation supplier: their actual
invariant categorical trace has the kernel(R09.14).
